\section{十一、项目模式}

本项目在技术路径与推广路径上分别形成了可复制的运作模式，二者共同构成从技术研发到基层服务的完整闭环。

\subsection{（1）技术研发模式：数据+模型双轮驱动}

\subsubsection{（1）以精准识别为核心，打造适配黑土场景的专属技术底座}

通用遥感变化检测模型多基于城市地表、道路、建筑等语义类别构建样本，其光谱与纹理统计特征与东北黑土区连片耕地差异显著：黑土耕地地表覆盖均质度高、作物物候引起的时相差异明显、地块边界依赖田埂与垄沟等弱特征，直接迁移通用模型易出现漏检与误检。为此，项目确立“数据先行、模型跟进、二者交替迭代”的研发思路，将数据建设与算法优化同步推进。

\textbf{数据支撑}：项目自建黑土影像数据集，针对东北黑土地域的光谱、纹理特征构建专属样本库，系统标注耕地撂荒、非农占用、侵蚀沟等典型变化类型，解决通用数据集对黑土区耕地变化识别不准的问题。样本库兼顾不同季节、不同作物与不同地形条件，使模型在真实业务场景中具备稳定的泛化能力；新采集的影像与基层核查反馈的变化图斑可回补入库，为模型持续迭代提供数据源头。

\textbf{算法优化}：项目采用改进的 BIT（Bipartite Transformer）双时相 Transformer 模型，以 ResNet18作为 CNN 骨干网络提取双时相影像的多尺度空间特征，并通过核心子模块“语义分词器”将特征图切分为语义 tokens，在前后两个时相之间建立长距离的语义关联，从而通过对比不同时期的卫星影像，精准捕捉耕地撂荒、非农占用等变化。该结构对地块内部的一致性与地块之间的边界变化均有较好响应，对细小、狭长的变化区域尤为敏感。在此基础上，项目持续迭代优化模型参数，在保证识别精确率达到既定指标的前提下，不断提升检测精度与效率，使模型在精度与算力开销之间取得更优平衡。

数据与模型的交替迭代构成研发模式的运转机制：新采集的影像经人工与算法协同标注进入样本库用于再训练，更新后的模型回流业务场景生成新的变化图斑，图斑经基层核查后再次回补入库。这一循环使精度提升不依赖一次性的大规模标注投入，而随监测范围扩大持续累积。项目选择轻量化骨干网络配合语义分词的组合，在识别精确率达标的前提下压缩算力开销，使模型兼顾专业核查与基层部署两类场景。

\subsection{（2）落地服务模式：平台+基层打通最后一公里}

技术成果只有进入基层工作流程才能转化为保护效益；现实中遥感监测成果往往停留在专业软件与专题图件中，形成“有发现、无响应”的断层。项目以轻量化、易使用为目标，从工具与流程两端降低使用门槛。

\textbf{（1）轻量化部署}：项目采用前后端分离架构，后端基于 FastAPI 构建服务接口，前端为可视化 Web 系统，功能覆盖影像上传、变化检测、结果可视化、AI 解读与一键报告生成，形成完整的业务操作链路。系统降低平台运行门槛，普通电脑即可流畅使用，无需专业设备支撑，适配乡镇、合作社等基层场景，使基层工作人员无需掌握遥感专业软件即可完成日常监测任务；一键报告功能将检测结果自动整理为规范文档，减少人工汇总与制表的重复劳动。

\textbf{（2）闭环预警服务}：项目向农户、合作社直接推送耕地撂荒、违规占用等预警信息，打破“监测—响应”的信息壁垒，打通耕地保护的“最后一公里”。预警信息并非孤立推送，而是与变化图斑的位置、面积、类型及处置建议一并送达，使基层主体能够直接开展核查与处置，形成发现、核实、反馈的完整链条，将原本分散的“监测”与“响应”连接为可追溯的责任链。

从运作流程看，落地服务模式由“监测—推送—核查—反馈”四个环节构成：系统按周更新影像并完成变化检测，将疑似变化图斑按行政区划与经营主体归集，推送至乡镇、合作社与农户；基层主体收到信息后开展实地核查，核查结论回填系统并作为后续监测与样本更新的依据。为保障服务可持续，项目以培训本地人员、扶持本地化服务团队的方式承接日常运维，使服务能力沉淀在基层，避免项目结束、服务中断。

\subsection{（3）项目核心支撑亮点}

两大模式的落地，依托七大核心特点形成完整支撑体系：

\begin{enumerate}
  \item \textbf{数据驱动}：以黑土专属影像数据为基础，为模型迭代和监测精度提供持续保障。
  \item \textbf{轻量部署}：算法与平台轻量化改造，降低基层使用门槛，实现“开箱即用”。
  \item \textbf{周级监测}：以周为周期更新影像数据，实现耕地变化的高频动态监测。
  \item \textbf{基层适配}：功能设计贴合乡镇、村集体的实际工作场景，操作简单易上手。
  \item \textbf{预警闭环}：从异常识别到信息推送，形成完整的监测预警与响应闭环。
  \item \textbf{国产遥感}：依托国产卫星影像数据，保障数据安全与自主可控。项目以国产高分二号（GF-2）亚米级光学遥感卫星影像为主要数据源，并辅以无人机航拍影像作为机动补充，在数据获取环节即实现自主可控。
  \item \textbf{用养结合}：将监测预警与耕地养护、政策服务相结合，实现“监测即服务”，助力黑土资源可持续保护。
\end{enumerate}

七大亮点构成层层递进的支撑逻辑：数据驱动与国产遥感回答“数据从哪里来”，轻量部署与基层适配回答“谁来用、是否用得起”，周级监测与预警闭环回答“发现之后怎么办”，用养结合回答“监测的最终目的”。

进一步看，七个亮点是按“数据基础—使用门槛—响应机制—最终目标”四个层次组织的：数据驱动与国产遥感位于底层，决定监测能力能否自主稳定供给；轻量部署与基层适配位于工具层，决定技术能否被一线人员真正掌握；周级监测与预警闭环位于机制层，决定发现的问题能否转化为实际处置；用养结合位于目标层，决定监测的落点是否回到黑土地质量本身。任一层次缺失，其他层次的成效都会被削弱。

\section{十二、项目核心收益：黑土智能检测系统}

本项目自研的黑土智能检测系统，通过 BIT 模型与轻量化部署技术，实现了黑土地变化监测的效率、成本与精度的三重突破，核心成果与收益如下：

\subsection{1.核心性能指标}

\begin{itemize}
  \item \textbf{识别精度}：BIT 模型对黑土地变化类别的识别精确率达70\%，可有效识别耕地撂荒、非农占用、侵蚀沟等关键变化类型，满足黑土地精细化监测的需求。该指标面向黑土区真实业务场景下的多类别变化识别任务，其意义在于系统输出的变化图斑中相当一部分可直接用于基层核查，从整体上降低人工复核工作量。
  \item \textbf{成本优化}：相较传统深度学习方案，算力成本降低约69\%，大幅降低了监测系统的部署与运行门槛。成本下降主要来自骨干网络的轻量化设计与前后端分离的轻量架构：前者减少单次推理的算力开销，后者使系统无需依赖高端硬件即可运行，从而把监测能力从专业机房延伸到基层办公环境。
  \item \textbf{效率提升}：从传统人工/常规遥感方法到本系统，单次监测周期缩短约33\%，实现了黑土地变化的高频动态监测。其价值不仅在于时间本身，更在于使“周级”更新成为常态化能力：监测频次提高意味着变化能够被更早发现，处置窗口相应前移，使耕地保护由事后统计转向事中干预。
\end{itemize}

三项指标并非彼此孤立，而是共同决定系统能否在真实业务中成立：识别精确率决定输出结果能否作为核查依据，算力成本决定系统能否下沉到基层办公环境，监测周期决定变化能否在处置窗口内被发现。三者结合，意味着系统在“看得准、用得起、发现得早”三方面同时具备可用性。

\subsection{2.验证与落地基础}

系统已在松嫩平原完成实地验证，适配普通电脑即可部署运行，无需高端硬件支撑，为后续在东北黑土区的规模化推广奠定了坚实基础。松嫩平原黑土分布集中、耕地经营主体类型丰富，其验证结论对东北其他黑土区具有较强的参考意义与可迁移性。

\subsection{3.业务与社会价值}

系统直接服务于黑土地保护“总量不减少、功能不退化”的核心目标，为黑土地常态化监测、耕地保护政策落地提供了智能化支撑，助力实现黑土地退化问题的“早发现、早预警、早治理”，兼具显著的生态效益与社会效益。

在业务层面，系统将《中华人民共和国黑土地保护法》《国家黑土地保护工程实施方案（2021—2025年）》《东北黑土地保护性耕作行动计划（2020—2025年）》等政策文件提出的监测与保护要求，转化为可操作、可查询、可追溯的信息化工具，使政策要求与基层执行之间有了明确的技术衔接点。在社会层面，系统面向农户与合作社免费提供预警服务，降低了小规模经营主体获取专业监测信息的门槛，使黑土地保护由少数专业机构的职责扩展为多方共同参与的日常实践。

\section{十三、项目成效与发展规划}

\subsection{1.当前项目成效}

赋能基层，带动多方价值。项目落地以来，已在农业帮扶与带动就业两方面取得了显著成效：

\begin{enumerate}
  \item \textbf{农业帮扶}：免费为农户推送撂荒、违法占用预警，已帮助500+农户提前止损；同时为10万亩黑土地提供“监测—诊断—措施”的全链条服务，助力黑土地用养结合，提升耕地质量。所谓“提前止损”，是指农户在违规行为造成不可逆后果之前即获得提示，有机会及时恢复耕种或纠正用地行为；全链条服务则使监测结果不止于发现问题，还能对接养护建议与政策服务。
  \item \textbf{带动就业}：通过培训基层农技员与返乡青年，新增智能监测岗位30+个；推广轻量化设备操作，扶持大学生创业与本地化服务团队；建立黑土地影像标注协作点，带动当地农户家门口就业，实现技术赋能与乡村振兴的结合。受训的基层人员既服务于本项目，也逐步成为当地黑土地保护工作的长期技术力量。
\end{enumerate}

农业帮扶与带动就业互为支撑：帮扶扩大监测服务覆盖面，岗位增设保障服务持续下沉到田间地头。

从可持续性看，两项成效共同指向“输血”与“造血”的结合：免费预警服务直接降低农户的损失风险，属于短期可感知的帮扶；岗位增设、设备操作培训与影像标注协作点则把监测工作拆解为基层可承接的环节，使农户与返乡青年在参与中获得技能与收入来源。影像标注协作点尤其关键，它既是模型迭代所需样本的生产端，也是本地就业的容纳端。

\subsection{2.项目发展历程与未来规划}

项目从2022年启动至今，已完成从调研到试点的关键阶段，并规划了清晰的长期发展路径：

\begin{table}[htbp]
  \centering
  \caption{项目发展历程与未来规划}
  \begin{tabular}{|c|p{0.68\textwidth}|}
  \hline
  时间节点 & 阶段目标与主要任务 \\
  \hline
  2022年 & 启动政策调研与数据积累，完成黑土地现状摸底 \\
  \hline
  2024年 & 完成影像标注与模型验证，实现技术闭环 \\
  \hline
  2026年 & 获得软件著作权并落地试点，形成成熟的产品形态 \\
  \hline
  2028年 & 推广至黑、吉、辽、蒙四省区，实现常态化监测 \\
  \hline
  2030年 & 构建“天空地”一体化监测网，实现东北黑土地全网动态监测 \\
  \hline
  \end{tabular}
  \label{tab:plan}
\end{table}

五个时间节点层层递进：2022年解决“有无依据与数据”，2024年解决“技术是否成立”，2026年解决“成果能否落地”，2028年解决“能否跨区域复制”，2030年则指向“能否形成全域覆盖的一体化监测能力”，前一阶段的产出构成后一阶段的输入。

具体来看，2022年的调研与数据积累为后续工作确立了地域范围与问题清单；2024年的影像标注与模型验证解决了技术是否成立的问题，使研发转为可复现的成果；2026年团队获软件著作权3项并落地试点，标志成果具备可交付的产品形态；2028年推广至黑、吉、辽、蒙四省区，重点在于把验证过的方案复制到气候、地形与经营方式不同的区域；2030年构建“天空地”一体化监测网，形成覆盖东北黑土区的常态化监测能力。

展望未来，项目将继续以双轮驱动夯实技术底座、以平台与基层贯通拓展应用边界，围绕黑土地保护“总量不减少、功能不退化”的核心目标持续发力，为东北黑土区的耕地保护提供长期、稳定、自主可控的技术支撑。
